\documentclass[a4paper]{article}

\def\npart{II}

\def\ntitle{Galois Theory}
\def\nlecturer{C.\ Brookes}

\def\nterm{Michaelmas}
\def\nyear{2017}

\ifx \nauthor\undefined
  \def\nauthor{Qiangru Kuang}
\else
\fi

\ifx \ntitle\undefined
  \def\ntitle{Template}
\else
\fi

\ifx \nauthoremail\undefined
  \def\nauthoremail{qk206@cam.ac.uk}
\else
\fi

\ifx \ndate\undefined
  \def\ndate{\today}
\else
\fi

\title{\ntitle}
\author{\nauthor}
\date{\ndate}

%\usepackage{microtype}
\usepackage{mathtools}
\usepackage{amsthm}
\usepackage{stmaryrd}%symbols used so far: \mapsfrom
\usepackage{empheq}
\usepackage{amssymb}
\let\mathbbalt\mathbb
\let\pitchforkold\pitchfork
\usepackage{unicode-math}
\let\mathbb\mathbbalt%reset to original \mathbb
\let\pitchfork\pitchforkold

\usepackage{imakeidx}
\makeindex[intoc]

%to address the problem that Latin modern doesn't have unicode support for setminus
%https://tex.stackexchange.com/a/55205/26707
\AtBeginDocument{\renewcommand*{\setminus}{\mathbin{\backslash}}}
\AtBeginDocument{\renewcommand*{\models}{\vDash}}%for \vDash is same size as \vdash but orginal \models is larger
\AtBeginDocument{\let\Re\relax}
\AtBeginDocument{\let\Im\relax}
\AtBeginDocument{\DeclareMathOperator{\Re}{Re}}
\AtBeginDocument{\DeclareMathOperator{\Im}{Im}}
\AtBeginDocument{\let\div\relax}
\AtBeginDocument{\DeclareMathOperator{\div}{div}}

\usepackage{tikz}
\usetikzlibrary{automata,positioning}
\usepackage{pgfplots}
%some preset styles
\pgfplotsset{compat=1.15}
\pgfplotsset{centre/.append style={axis x line=middle, axis y line=middle, xlabel={$x$}, ylabel={$y$}, axis equal}}
\usepackage{tikz-cd}
\usepackage{graphicx}
\usepackage{newunicodechar}

\usepackage{fancyhdr}

\fancypagestyle{mypagestyle}{
    \fancyhf{}
    \lhead{\emph{\nouppercase{\leftmark}}}
    \rhead{}
    \cfoot{\thepage}
}
\pagestyle{mypagestyle}

\usepackage{titlesec}
\newcommand{\sectionbreak}{\clearpage} % clear page after each section
\usepackage[perpage]{footmisc}
\usepackage{blindtext}

%\reallywidehat
%https://tex.stackexchange.com/a/101136/26707
\usepackage{scalerel,stackengine}
\stackMath
\newcommand\reallywidehat[1]{%
\savestack{\tmpbox}{\stretchto{%
  \scaleto{%
    \scalerel*[\widthof{\ensuremath{#1}}]{\kern-.6pt\bigwedge\kern-.6pt}%
    {\rule[-\textheight/2]{1ex}{\textheight}}%WIDTH-LIMITED BIG WEDGE
  }{\textheight}% 
}{0.5ex}}%
\stackon[1pt]{#1}{\tmpbox}%
}

%\usepackage{braket}
\usepackage{thmtools}%restate theorem
\usepackage{hyperref}

% https://en.wikibooks.org/wiki/LaTeX/Hyperlinks
\hypersetup{
    %bookmarks=true,
    unicode=true,
    pdftitle={\ntitle},
    pdfauthor={\nauthor},
    pdfsubject={Mathematics},
    pdfcreator={\nauthor},
    pdfproducer={\nauthor},
    pdfkeywords={math maths \ntitle},
    colorlinks=true,
    linkcolor={red!50!black},
    citecolor={blue!50!black},
    urlcolor={blue!80!black}
}

\usepackage{cleveref}



% TODO: mdframed often gives bad breaks that cause empty lines. Would like to switch to tcolorbox.
% The current workaround is to set innerbottommargin=0pt.

%\usepackage[theorems]{tcolorbox}





\usepackage[framemethod=tikz]{mdframed}
\mdfdefinestyle{leftbar}{
  %nobreak=true, %dirty hack
  linewidth=1.5pt,
  linecolor=gray,
  hidealllines=true,
  leftline=true,
  leftmargin=0pt,
  innerleftmargin=5pt,
  innerrightmargin=10pt,
  innertopmargin=-5pt,
  % innerbottommargin=5pt, % original
  innerbottommargin=0pt, % temporary hack 
}
%\newmdtheoremenv[style=leftbar]{theorem}{Theorem}[section]
%\newmdtheoremenv[style=leftbar]{proposition}[theorem]{proposition}
%\newmdtheoremenv[style=leftbar]{lemma}[theorem]{Lemma}
%\newmdtheoremenv[style=leftbar]{corollary}[theorem]{corollary}

\newtheorem{theorem}{Theorem}[section]
\newtheorem{proposition}[theorem]{Proposition}
\newtheorem{lemma}[theorem]{Lemma}
\newtheorem{corollary}[theorem]{Corollary}
\newtheorem{axiom}[theorem]{Axiom}
\newtheorem*{axiom*}{Axiom}

\surroundwithmdframed[style=leftbar]{theorem}
\surroundwithmdframed[style=leftbar]{proposition}
\surroundwithmdframed[style=leftbar]{lemma}
\surroundwithmdframed[style=leftbar]{corollary}
\surroundwithmdframed[style=leftbar]{axiom}
\surroundwithmdframed[style=leftbar]{axiom*}

\theoremstyle{definition}

\newtheorem*{definition}{Definition}
\surroundwithmdframed[style=leftbar]{definition}

\newtheorem*{slogan}{Slogan}
\newtheorem*{eg}{Example}
\newtheorem*{ex}{Exercise}
\newtheorem*{remark}{Remark}
\newtheorem*{notation}{Notation}
\newtheorem*{convention}{Convention}
\newtheorem*{assumption}{Assumption}
\newtheorem*{question}{Question}
\newtheorem*{answer}{Answer}
\newtheorem*{note}{Note}
\newtheorem*{application}{Application}

%operator macros

%basic
\DeclareMathOperator{\lcm}{lcm}

%matrix
\DeclareMathOperator{\tr}{tr}
\DeclareMathOperator{\Tr}{Tr}
\DeclareMathOperator{\adj}{adj}

%algebra
\DeclareMathOperator{\Hom}{Hom}
\DeclareMathOperator{\End}{End}
\DeclareMathOperator{\id}{id}
\DeclareMathOperator{\im}{im}
\DeclarePairedDelimiter{\generation}{\langle}{\rangle}

%groups
\DeclareMathOperator{\sym}{Sym}
\DeclareMathOperator{\sgn}{sgn}
\DeclareMathOperator{\inn}{Inn}
\DeclareMathOperator{\aut}{Aut}
\DeclareMathOperator{\GL}{GL}
\DeclareMathOperator{\SL}{SL}
\DeclareMathOperator{\PGL}{PGL}
\DeclareMathOperator{\PSL}{PSL}
\DeclareMathOperator{\SU}{SU}
\DeclareMathOperator{\UU}{U}
\DeclareMathOperator{\SO}{SO}
\DeclareMathOperator{\OO}{O}
\DeclareMathOperator{\PSU}{PSU}

%hyperbolic
\DeclareMathOperator{\sech}{sech}

%field, galois heory
\DeclareMathOperator{\ch}{ch}
\DeclareMathOperator{\gal}{Gal}
\DeclareMathOperator{\emb}{Emb}



%ceiling and floor
%https://tex.stackexchange.com/a/118217/26707
\DeclarePairedDelimiter\ceil{\lceil}{\rceil}
\DeclarePairedDelimiter\floor{\lfloor}{\rfloor}


\DeclarePairedDelimiter{\innerproduct}{\langle}{\rangle}

%\DeclarePairedDelimiterX{\norm}[1]{\lVert}{\rVert}{#1}
\DeclarePairedDelimiter{\norm}{\lVert}{\rVert}



%Dirac notation
%TODO: rewrite for variable number of arguments
\DeclarePairedDelimiterX{\braket}[2]{\langle}{\rangle}{#1 \delimsize\vert #2}
\DeclarePairedDelimiterX{\braketthree}[3]{\langle}{\rangle}{#1 \delimsize\vert #2 \delimsize\vert #3}

\DeclarePairedDelimiter{\bra}{\langle}{\rvert}
\DeclarePairedDelimiter{\ket}{\lvert}{\rangle}




%macros

%general

%divide, not divide
\newcommand*{\divides}{\mid}
\newcommand*{\ndivides}{\nmid}
%vector, i.e. mathbf
%https://tex.stackexchange.com/a/45746/26707
\newcommand*{\V}[1]{{\ensuremath{\symbf{#1}}}}
%closure
\newcommand*{\cl}[1]{\overline{#1}}
%conjugate
\newcommand*{\conj}[1]{\overline{#1}}
%set complement
\newcommand*{\stcomp}[1]{\overline{#1}}
\newcommand*{\compose}{\circ}
\newcommand*{\nto}{\nrightarrow}
\newcommand*{\p}{\partial}
%embed
\newcommand*{\embed}{\hookrightarrow}
%surjection
\newcommand*{\surj}{\twoheadrightarrow}
%power set
\newcommand*{\powerset}{\mathcal{P}}

%matrix
\newcommand*{\matrixring}{\mathcal{M}}

%groups
\newcommand*{\normal}{\trianglelefteq}
%rings
\newcommand*{\ideal}{\trianglelefteq}

%fields
\renewcommand*{\C}{{\mathbb{C}}}
\newcommand*{\R}{{\mathbb{R}}}
\newcommand*{\Q}{{\mathbb{Q}}}
\newcommand*{\Z}{{\mathbb{Z}}}
\newcommand*{\N}{{\mathbb{N}}}
\newcommand*{\F}{{\mathbb{F}}}
%not really but I think this belongs here
\newcommand*{\A}{{\mathbb{A}}}

%asymptotic
\newcommand*{\bigO}{O}
\newcommand*{\smallo}{o}

%probability
\newcommand*{\prob}{\mathbb{P}}
\newcommand*{\E}{\mathbb{E}}

%vector calculus
\newcommand*{\gradient}{\V \nabla}
\newcommand*{\divergence}{\gradient \cdot}
\newcommand*{\curl}{\gradient \cdot}

%logic
\newcommand*{\yields}{\vdash}
\newcommand*{\nyields}{\nvdash}

%differential geometry
\renewcommand*{\H}{\mathbb{H}}
\newcommand*{\transversal}{\pitchfork}
\renewcommand{\d}{\mathrm{d}} % exterior derivative

%number theory
\newcommand*{\legendre}[2]{\genfrac{(}{)}{}{}{#1}{#2}}%Legendre symbol


\DeclareMathOperator{\n}{N}

\begin{document}

\begin{titlepage}
  \begin{center}
    \includegraphics[width=0.6\textwidth]{logo.jpg}\par
    \vspace{1cm}
    {\scshape\huge Mathematics Tripos \par}
    \vspace{2cm}
    {\huge Part \npart \par}
    \vspace{0.6cm}
    {\Huge \bfseries \ntitle \par}
    \vspace{1.2cm}
    {\Large\nterm, \nyear \par}
    \vspace{2cm}
    
    {\large \emph{Lectures by } \par}
    \vspace{0.2cm}
    {\Large \scshape \nlecturer}
    
    \vspace{0.5cm}
    {\large \emph{Notes by }\par}
    \vspace{0.2cm}
    {\Large \scshape \href{mailto:\nauthoremail}{\nauthor}}
 \end{center}
\end{titlepage}

\tableofcontents

\section{Field Extensions}

\blindtext{10}

\section{Seperable, Normal and Galois Extensions}

\blindtext

\subsection{Trace \& Norm}

\begin{definition}[Trace and norm]
  Let \(K \leq M\) be a finite field extension and \(\alpha \in M\). Multiplication by \(a\) is a \(K\)-linear map \(\theta_\alpha: M \to M\). The \emph{trace of \(\alpha\) over \(K\)} is
  \[
    \Tr_{M/K}(\alpha) = \Tr( \theta_\alpha) \in K
  \]
  and the \emph{norm of \(\alpha\) over \(K\)} is
  \[
    \n_{M/K}(\alpha) = \det \theta_\alpha \in K.
  \]
\end{definition}

\begin{note}
  These dependent on the field extension.
\end{note}

\begin{theorem}
  Suppose \(f_\alpha(t) = t^s + a_{s - 1}t^{s - 1} + \dots + a_0\) is the minimal polynomial of \(\alpha\) over \(K\). Let \(r = |M: K(\alpha)|\). Then the characteristic polynomial of \(\theta_\alpha\) is \((f_\alpha(t))^r\) and
  \begin{align*}
    \Tr_{M/K}(\alpha) &= -r a_{s-1} \\
    \n_{M/K}(\alpha) &= ((-1)^s a_0)^r
  \end{align*}
\end{theorem}

\subsection{Normal Extensions}

We met this definition before:

\begin{definition}[Normal extension]
  An extension \(K \leq M\) is \emph{normal} if for every \(\alpha \in M\), the minimal polynomial \(f_\alpha(t)\) of \(\alpha\) over \(K\) splits over \(M\).
\end{definition}

\begin{theorem}
  Let \(K \leq M\) be a finite field extension. Then \(K \leq M\) is normal if and only if \(M\) is the splitting field for some \(f(t) \in K[t]\).
\end{theorem}

\begin{proof}\leavevmode
  \begin{itemize}
  \item \(\Rightarrow\): suppose \(K \leq M\) isnormal. Pick \(\alpha_1, \dots, \alpha_r \in M\) such that \(M = K(\alpha_1, \dots, \alpha_kr\). Let \(f(_{\alpha_i}(t)\) be the minimal polynomail of \(\alpha_i\) over \(K\). Let \(f(t) = \prod_{n = 1}^{r} f_{\alpha_i}(t)\).

    By normality, each \(f_{\alpha_i}(t)\) splits over \(M\) and so is \(f(t)\). \(M\) is the splitting field of \(f(t)\) over \(K\) since if \(\beta_1, \dots, \beta_m\) are the roots of \(f(t)\) then \(M = K(\beta_1, \dots, \beta_m)\).
  \item \(\Leftarrow\): suppose \(M\) is a splitting field for \(f(t)\) over \(K\). Then \(M = K(\beta_1, \dots, \beta_m)\) where \(\beta_j\) are the roots of \(f(t)\) over \(M\). Take \(\alpha \in M\). Let \(f_\alpha(t)\) be the minimal polynomial of \(\alpha\) over \(K\). Let \(M \leq L\) large enough so that \(f_\alpha(t)\) splits over \(L\)

    Now consdier a \(K\)-homomorphism \(\phi: M \to L\): \(\phi(\beta_j)\) is also a root of \(f(t)\) and is therefore one of the \(\beta_j\). Injectivity of \(K\)-homomorphisms implies that \(\phi\) permutes \(\beta_j\). However, \(M = K(\beta_1, \dots, \beta_m)\) and so \(\phi\) is determined by the images of \(\beta_j\), thus \(\phi(M) = M\).

    However, if \(\alpha_i\) is a root of \(f_\alpha(t)\) in \(L\), then there is a \(K\)-homomorphism
    \begin{align*}
      K(\alpha) &\to K(\alpha_i) \\
      \alpha &\mapsto \alpha_i
    \end{align*}
    This extends by 2.10~\ref{tbf} to a \(K\)-homomorphism \(\phi: M \to L\). But \(\phi(M) = M\) so \(\alpha_i \in M\). Thus \(M\) is normal over \(K\).
  \end{itemize}
\end{proof}

\begin{remark}
  As for separability, being normal is equivalent to being normally generated. We will show in example sheet that \(K \leq L\) is a normal and finite extension if and only if \(L = K(\alpha_1, \dots, \alpha_r)\) with the minimal polynomial of each adjoined element splitting over \(L\).
\end{remark}

\begin{definition}[\(K\)-automorphism group]
  Let \(K \leq M\) be a finite extension. Its \emph{\(K\)-automorphism group} is
  \[
    \aut_K(M) = \Hom_K(M, M).
  \]
\end{definition}

From \ref{tbh} we know that such \(K\)-automorphisms are isormophisms and thus have inverses (well, it name says so).

\begin{lemma}
  \[
    |\aut_K(M)| \leq |M:K|.
  \]
\end{lemma}

\begin{proof}
  \ref{tbh} (extension of homomorphisms)
\end{proof}

\begin{theorem}
  \label{thm:galois criterion}
  Let \(K \leq M\) be a finite field extension, then
  \[
    |\aut_K(M)| = |M:K|
  \]
  if and only if the extension is both normal and separable.
\end{theorem}

\begin{definition}[Galois extension]
  A finite field extension that is normal and separable is a \emph{Galois extension}.
\end{definition}

\begin{definition}[Galois group]
  Let \(K \leq M\) be a Galois extension. Then the \(K\)-automorphism group of \(M\) is the \emph{Galois group} of \(M\) over \(K\), denoted by
  \[
    \gal(M/K).
  \]
\end{definition}

\begin{remark}
  Some authors use the term ``Galois group'' as a synonym for automorphism group even when the extension is not Galois.
\end{remark}

\begin{proof}[Proof of \Cref{thm:galois criterion}]
  Suppose \(|\aut_K(M)| = |M:K| = n\). Let \(M \leq L\) be large enough. The \(n\) distinct \(K\)-automorphisms \(\phi: M \to M\) extend to \(n\) \(K\)-homomorphisms \(\phi: M \to L\) and 2.12~\ref{tbf} says that \(M\) is seperable over \(K\).

  For normality, pick \(\alpha \in M\) with minimal polynomial \(f_\alpha(t)\) over \(K\). Take \(M = K(\alpha_1, \dots, \alpha_m)\) as in the proof of 2.10~\ref{tbf} with \(\alpha = \alpha_1\) and \(L = M\). We only get \(|M:K|\) extensions of the inclusion \(K \hookrightarrow M\) if each inequality in the proof is an equality. In particular, we need the number of \(K\)-homomorphisms \(K(\alpha_1) \to M\) to be \(K(\alpha_1):K|\). But then 2.6~\ref{tbf} says we have \(|K(\alpha):K|\) distinct roots of \(f_\alpha(t)\) i \(M\). Thus \(f_\alpha(t)\) splits over \(M\).

  Conversely, suppose \(K \leq M\) is separable and normal. Then for \(K \leq M \leq L\) with \(L\) large enough, separability implies there are \(|M:K|\) \(K\)-homomorphisms \(\phi: M \to L\) by 2.12~\ref{tbf}. However, \(K \leq M\) is normal implies that it is the splitting field for some polynomial \(f(t) \in K[t]\) and thus \(M = K(\alpha_1, \dots, \alpha_n)\) where \(f(t) = \prod_{i = 1}^n (t - \alpha_i)\). Note that \(\phi(\alpha_j)\) is also a root of \(\phi(f(t)) = f(t)\), and is therefore one of the \(\alpha_j\). Thus \(\phi(M) = M\). Thus \(|\aut_K(M)| = |M:K|\).
\end{proof}

\begin{remark}
  In the previous proof we have shown that if \(K \leq M \leq L\), \(\phi \in \Hom_K(M, L)\) and \(K \leq M\) is normal then \(\phi(M) = M\).
\end{remark}

\begin{eg}\leavevmode
  \begin{enumerate}
  \item Consider \(\Q \leq \Q(\sqrt 2, i)\), which is Galois:
    \[
      \gal(\Q(\sqrt 2, i)/\Q) = \generation{\sigma: \sqrt 2 \mapsto -\sqrt 2, \tau: i \mapsto -1} \cong C_2 \times C_2.
    \]
    All non-identity elements have order \(2\).
  \item Let \(f(t) = t^3 - 2\). The splitting field of \(f(t)\) over \(\Q\) is \(\Q(\sqrt[3]{2}, \omega)\) where \(\omega\) is a primitive cubic root of unity. Thus \(\Q \leq \Q(\sqrt[3]{2}, \omega)\) is Galois with \(|\gal(\Q(\sqrt[3]{2}, \omega)/\Q) = |\Q(\sqrt[3]{2}, \omega):\Q| = 6\). The Galois group contains
    \begin{align*}
      \sigma_1: \sqrt[3]{2} \mapsto \sqrt[3]{2} %%%%%%%%
    \end{align*}
    \texttt{to be filled in}
    %
    %
    %
    %
    %
    %
  \end{enumerate}
\end{eg}
\section{Fundamental Theorem of Galois Theory}

\begin{definition}[Fixed field]
  Let \(K \leq L\) be a field extension and \(H \leq \aut_K(L)\). The \emph{fixed field} of \(H\) is
  \[
    L^H = \{\alpha \in L: \sigma(\alpha) = \alpha \text{ for all } \sigma \in H \}.
  \]
\end{definition}

The fixed field is a field and \(K \leq L^H \leq L\).

\begin{theorem}[Fundamental Theorem of Galois Theory]
  \label{thm:fundamental}
  Let \(K \leq L\) be a finite Galois extension. Then
  \begin{enumerate}
  \item There is a one-to-one correspondence
    \begin{align*}
      \left\{
        \begin{array}[h]{c}
          \text{intermediate subfield} \\
          K \leq M \leq L
        \end{array}
      \right\}
      &\longleftrightarrow
      \left\{
        \begin{array}[h]{c}
          \text{subgroup } H \\
          \text{of } \gal(L/K)
        \end{array}
      \right\} \\
      %
      M &\mapsto \aut_M(L) \\
      L^H &\mapsfrom H
    \end{align*}
  \item \(H\) is a normal subgroup of \(\gal(L/K)\) if and only if \(K \leq L^H\) is normal if and only if \(K \leq L^H\) is Galois.
  \item If \(H \trianglelefteq \gal(L/K)\) then the map
    \[
      \theta: \gal(L/K) \to \gal(L^H/K)
    \]
    given by restriction to \(L^H\) is a surjective group homomorphism with kernel \(H\).
  \end{enumerate}
\end{theorem}

\begin{remark}\leavevmode
  \begin{enumerate}
  \item Observe that \(M \leq L\) is Galois and so we could have written \(\gal(L/M)\) instead of \(\aut_M(L)\). To see this,
    \begin{itemize}
    \item seperability: follows from \ref{tbf}
    \item normality: if \(\alpha \in L\) then the minimal polynomial of \(\alpha\) over \(M\) divides the minimal polynomial of \(\alpha\) over \(K\). But the latter splits over \(L\) (see example sheet).
    \end{itemize}
  \item If \(K \leq M\) is normal then \ref{tbf} says if \(\sigma: L \to L\) then \(\sigma(M) = M\) and so we can talk about the restriction of \(\alpha\) to \(M\) giving an automorphism of \(M\).
  \end{enumerate}
\end{remark}

\begin{eg}\leavevmode
  \begin{enumerate}
  \item \(\Q \leq \Q(\sqrt{2}, i)\). We saw in \ref{tbf} the lattices of intermediate subfields and subgroups \(\gal(\Q(\sqrt 2, i/\Q) \cong C_2 \times C_2\) which is abelian. Thus all subgroups are normal and all intermediate subfields are normal extensions of \(\Q\).
  \item \(\Q \leq \Q(\sqrt[3]{2}, \omega)\).
    \[
      \begin{tikzcd}
        & \Q(\sqrt[3]{2}, \omega) \ar[dl, dash, "3"] \ar[d, dash, "2"] \ar[dr, dash, "2"] \ar[drr, dash, "2"] \\
        \Q(\omega) \ar[dr, dash, "2"] & \Q(\sqrt[3]{2}) \ar[d, dash, "3"] & \Q(\omega\sqrt[3]{2}) \ar[dl, dash, "3"] & \Q(\omega^2\sqrt[3]{2}) \ar[dll, dash, "3"] \\
        & \Q
      \end{tikzcd}
    \]
    \[
      \begin{tikzcd}
        & 1 \ar[dl, dash, "3"] \ar[d, dash, "2"] \ar[dr, dash, "2"] \ar[drr, dash, "2"] \\
        \{\sigma_1, \sigma_2, \sigma_3\} \ar[dr, dash, "2"] & \{\sigma_1, \sigma_4\} \ar[d, dash, "3"] & \{\sigma_1, \sigma_5\} \ar[dl, dash, "3"] & \{\sigma_1, \sigma_6\} \ar[dll, dash, "3"] \\
        & \{\sigma_1, \dots, \sigma_6\} \cong D_6
      \end{tikzcd}
    \]
    The subgroup \(H\) of order \(3\) is normal but those of order \(2\) are not so the map \(\gal(\Q(\sqrt[3]{2}, \omega)/\Q) \to \gal(\Q(\omega)/\Q)\) has kernel \(H\).
    \end{enumerate}
\end{eg}

\begin{theorem}[Artin's]
  \label{thm:artin}
  Let \(K \leq L\) be a field extension and \(H \leq \aut_K(L)\) be a finite subgroup. Let \(M = L^H\). Then \(M \leq L\) is a finite Galois extension and \(H = \gal(L/M)\).
\end{theorem}

\begin{remark}
  This implies some part of the Galois correspondence:
  \[
    H \to L^H \to \gal(L/H^H)
  \]
  where we get back to \(H\).
\end{remark}

\begin{proof}
  Take \(\alpha \in L\). We first show that \(|M(\alpha):M| \leq |H|\). Let \(\{\alpha_1, \dots, \alpha_n\}\) be the distinct images of \(\alpha\) under \(H\), i.e.\ \(\{\phi(\alpha): \alpha \in H\}\). Define \(g(t) = \prod_{i = 1}^{\infty} (t - \alpha_i) \). Each \(\phi\) induces an endomorphism on \(L[t]\) under which \(g(t)\) is invariant since \(\phi\) permutes the \(\alpha_i\). Thus the coefficients lie in \(L^H = M\) and so \(g(t) \in M[t]\). By definition \(g(\alpha) = 0\) since \(\alpha\) is one of the \(\alpha_i\). Hence the minimal polynomial \(f_\alpha(t)\) of \(\alpha\) over \(M\) divides \(g(t)\). Thus
  \[
    |M(\alpha):M| = \deg f_\alpha(t) \leq \deg g(t) \leq |H|.
  \]
  This step shows that \(\alpha\) is algebraic over \(M\) and \(f_\alpha(t)\) is seperable whenver \(g(t)\) is. Thus \((M \leq L\) is a separable extension.

  Next we show that \(M \leq L\) is a simple extension. Pick \(\alpha \in L\) with \(|M(\alpha):m|\) maximal. We will show that \(L = M(\alpha)\) for this \(\alpha\). Suppose \(\beta \in L\). Then \(M \leq M(a=\alpha, \beta)\) is finite and separably generated and hence is a finite separable extension. By \nameref{thm:primitive} \(M(\alpha, \beta) = M(\gamma)\) for some \(\gamma\). But then
  \[
    M \leq M(\alpha) \leq M(\gamma)
  \]
  so by the maximality of \(|M(\alpha):M|\), \(M(\alpha) = M(\gamma)\). Thus \(\beta \in M(\gamma) = M(\alpha)\) so \(L = M(\alpha)\).

  It follows that \(|L:M| \leq |H|\).

  Finally,
  \[
    |L:M| = |M(\alpha):M| \leq |H| \leq |\aut_M(L)| \leq |L:M| \leq |L:M|
  \]
  so we must have equality throughtout, i.e.\ \(|L:M| = |\aut_M(L)| = |H|\) so \(M \leq L\) is a finite Glois extension and \(H = \gal(L/M)\).
\end{proof}

\begin{theorem}
  Let \(K \leq L\) be a finite field extension. TFAE:
  \begin{enumerate}
  \item \(K \leq L\) is Galois,
  \item \(L^H = K\) when \(H = \aut_K(L)\).
  \end{enumerate}
\end{theorem}

\begin{remark}
  The theorem allows some authors to give yet another definition of a Galois extension.
\end{remark}

\begin{proof}\leavevmode
  \begin{enumerate}
  \item \(1 \Rightarrow 2\): Let \(M = L^H\) where \(H = \aut_K(L)\). By \nameref{thm:artin} \(M \leq L\) is a Galois extension. Now we have two Galois extensions, giving the equality
    \begin{align*}
      |H| &= |\gal_K(L)| = |L:K| \\
          &= |\gal_M(L)| = |L:M|
    \end{align*}
    As \(K \leq L^H = M\), equality.
  \item \(2 \Rightarrow 1\): \nameref{thm:artin}.
  \end{enumerate}
\end{proof}

\begin{proof}[Proof of \nameref{thm:fundamental}]\leavevmode
  \begin{enumerate}
  \item Composing the maps
    \begin{align*}
      H &\mapsto L^H \\
      M &\mapsto \gal(L/M)
    \end{align*}
    gives \(H \to H\) by \nameref{thm:artin}. Composition the other way \(M \mapsto \gal(L/M) \mapsto L^H\) where \(H = \gal(L/M)\) gives \(M\) since \(M \leq L^H\) and
    \[
      |L:L^H| = |H| = |\gal(L/M)| = |L:M|.
    \]
  \item Take \(H \leq \gal(L/K)\). Then for \(\phi \in \gal(L/K)\), \(L^{\phi H \phi^{-1}} = \phi(L^H)\) so by 1 \(H \normal \gal(L/K)\) if and only if \(\phi(L^H) = L^H\). Let \(M = L^H\). We will show that \(K \leq M\) is normal if and only if \(\phi(M) = M\) for every \(\phi \in \gal(L/K)\):
    \begin{itemize}
    \item \(\Rightarrow\): To be filled in
      %
      %
      %
      %
      %
      %
    \end{itemize}

  \end{enumerate}
\end{proof}

\section{Galois Group of Polynomials}

\begin{definition}[Galois group]
  Let \(f(t) \in K[t]\) be a separable polynomial and \(K \leq L\) with \(L\) a splitting field for \(f(t)\). Then the \emph{Galois group} of \(f(t)\) over \(K\) is
  \[
    \gal(f) = \gal(L/K).
  \]
\end{definition}

Since \(L\) is a splitting field for \(f(t)\), \(L = K(\alpha_1, \dots, \alpha_n)\) where \(\alpha_1, \dots, \alpha_n\) are the roots of \(f(t)\) in \(L\). Observe that if \(\phi \in \gal(L/K)\) it maps bijectively roots of \(f(t)\) to itself. Thus \(\phi\) permutes the \(\alpha_i\).

Moreover if \(\phi\) fixes each \(\alpha_i\) it also fixes all elements of \(L\) and so it is the identity map. Thus \(\gal(f)\) may be regarded as a permutation group of the \(n\) roots, so in particular admitting a permutation representation in \(S_n\).

\begin{lemma}
  \label{lem:galois orbits}
  Suppose seperable \(f(t) = g_1(t) \cdots g_s(t)\) with \(g_i(t)\) irreducible in \(K[t]\) is a factorisation in \(K[t]\). Then the orbits of \(\gal(f)\) acting on the roots of \(f(t)\) correspond to the factors \(g_j(t)\): two roots are in the same orbit if and only if they are roots of the same \(g_i(t)\).

  In particular if \(f(t) \in K[t]\) is irreducible, there is only one orbit, i.e.\ \(\gal(f)\) acts transitively on the roots of \(f(t)\).
\end{lemma}

\begin{proof}
  Let \(\alpha_k\) and \(\alpha_\ell\) be in the same orbit under \(\gal(f)\). Thus there is \(\phi \in \gal(f)\) with \(\alpha_\ell = \phi(\alpha_k)\). But if \(\alpha_k\) is a root of \(g_j(t)\) then \(\phi(\alpha_k)\) is also a root of \(g_j(t)\).

  Conversely if \(\alpha_k\) and \(\alpha_\ell\) are roots of \(g_j(t)\), then
  \[
    K(\alpha_k) \cong K[t]/(g_j(t)) \cong K(\alpha_\ell) \leq L
  \]
  Denote by \(\phi_0\) the isomorphsim \(K(\alpha_k) \to K(\alpha_\ell)\). \(\phi_0\) extends to \(\phi \in \gal(L/K)\). Thus \(\alpha_k\) and \(\alpha_\ell\) are in the same orbit.
\end{proof}

\begin{lemma}
  The transitive subgroups of \(S_n\) for \(n \leq 5\) are
  \[
  \begin{array}[h]{c|l}
    n & \\ \hline
    2 & S_2 \cong C_2 \\ \hline
    3 & A_3 \cong C_3, S_3 \\ \hline
    4 & C_4, V_4, D_8, A_4, S_4 \\ \hline
    5 & C_5, D_{10}, H_{20}, A_5, S_5 \\ \hline
  \end{array}
  \]
  where \(H_{20}\) is generated by a \(5\)-cycle and a \(4\)-cycle.
\end{lemma}

\begin{theorem}
  Let \(p\) be a prime and \(f(t) \in \Q[t]\) of degree \(p\). Suppose \(f(t)\) has exactly \(2\) non-real roots in \(\C\). Then
  \[
    \gal(f) \cong S_p.
  \]
\end{theorem}

\begin{proof}
  \(\gal(f)\) acts on the \(p\) distincts roots of \(f(t)\) in a splitting field \(L\) of \(f(t)\) (in \(\C\)). By \Cref{lem:galois orbits} the irreducibility of \(f(t)\) implies that \(\gal(f)\) acts transitively on \(p\) roots so by orbit-stabiliser \(p \divides |\gal(f)|\) but
  \[
    |\gal(f)| \leq |S_p| = p!
  \]
  and so \(\gal(f)\) has a Sylow \(p\)-subgroup of order \(p\), necessarily cyclic. Thus \(\gal(f)\) contains a \(p\)-cycle. Supposition that there are exactly \(2\) non-real roots gives that complex conjugation yields a transposition in \(\gal(f)\). The \(p\)-cycle and the transposition generate \(S_p\).
\end{proof}

\begin{eg}
  Let \(f(t) = t^5 - 6t + 3 \in \Q[t]\), then claim \(\gal(f) \cong S_5\):

  \begin{proof}
    \(f(t)\) is irreducible by Einsenstein with \(p = 3\). We want to show that \(f(t)\) has \(3\) real roots (so \(2\) non-real roots) and apply the previous theorem.

    Now we apply knowledge from analysis:
    \[
      f(-2) = -17, f(-1) = 8, f(1) = -2, f(2) = 23
    \]
    and \(f'(t) = 5t^4 - 6\) which has two real roots. Thus by Intermediate Value Theorem there are \(3\) real roots while by Rolle's Theorem there are at most \(3\) real roots.
  \end{proof}
\end{eg}

\begin{definition}[Discriminant]
  Let \(f(t) \in K[t]\) with distinct roots \(\alpha_1, \dots, \alpha_n\) (in a splitting field) (\(f(t)\) need not be irreducible). Set \(\Delta = \prod_{i < j}(\alpha_i - \alpha_j)\). The the \emph{discriminant} \(D = D(f)\) of \(f\) is
  \[
    D(f) = \Delta^2 = \prod_{i < j}(\alpha_i - \alpha_j)^2 = (-1)^{\frac{n(n-1)}{2}} \prod_{i \neq j}(\alpha_i - \alpha_j).
  \]
\end{definition}

\begin{remark}
  We have already met this in the proof of~\ref{tbf}.
\end{remark}

\begin{lemma}
  Let \(f(t) \in K[t]\) be seperable of degree \(n\) with \(\ch K \neq 2\). Then
  \[
    \gal(f) \leq A_n \Leftrightarrow D(f) \text{ is a square in } K.
  \]
\end{lemma}

\begin{proof}
  Let \(L\) be a splitting field of \(f(t)\) over \(K\). Then \(D(f) \neq 0\) and is fixed by all elements of \(G = \gal(L/K)\) as the latter permutes the roots. Thus \(D \in K\) since \(L^G = K\) by Galois correspondence.

  If \(\sigma \in G\) then \(\sigma(\Delta) = (\sgn \sigma) \Delta\) where we are regarding \(G\) as a subgroup of \(S_n\) and \(\sgn\) is the signature (this is where we need \(\ch K \neq 2\)). Thus if \(G \leq A_n\) we got that \(\Delta\) is fixed by all \(\sigma \in G\). THus \(\Delta \in K = L^G\).

  On the other hand if \(G \nleq A_n\) we get \(\sigma(\Delta) = -\Delta\) if \(\sigma\) is odd and so \(\Delta \notin K = L^G\). Finally note that if \(D\) has square roots they must be \(\pm \Delta\).
\end{proof}

\begin{eg}\leavevmode
  \begin{enumerate}
  \item \(n = 2\): \(f(t) =t^2 + bt + c = (t - \alpha_1)(t - \alpha_2)\) has determinant
    \[
      D(f) = (\alpha_1 - \alpha_2)^2 - (\alpha_1 + \alpha_2)^2 - 4\alpha_1\alpha_2 = b^2 - 4c
    \]
  \item \(n = 3\): \(f(t) = t^3 + ct + d\) has determinant
    \[
      D(t) = -4c^3 - 27d^2
   \]
   \begin{remark}
     Any general monic cubic \(g(t)\) can be put into this form by a suitable substitution \(f(t) = g(t_1 + \lambda)\) for suitable \(\lambda\). Note \(D(f) = D(g)\).
   \end{remark}
 \end{enumerate}
\end{eg}

\begin{eg}\leavevmode
  \begin{itemize}
  \item \(f(t) = t^3 - t - 1 \in \Q[t]\) irreducible in \(\Z[t]\) since irreducible mod \(2\). \(D(f) = -23\) which is not a square in \(\Q\). Thus \(\gal(f) \cong S_3\).
  \item \(f(t) = t^3 - 3t - 1 \in \Q[t]\) irredicible since irreducible mod \(2\). \(D(f) = 81\) which is a square so \(\gal(f) \cong A_3\).
  \end{itemize}
\end{eg}



\end{document}